% 第3章 システム。oshifit_spec.md 第3節にもとづく下書き。
\chapter{OshiFit の設計}
\section{全体の流れ}
OshiFit は、最初の調査、アバターの作成、分析画面、2つの未来の提示、振り返り、最後の調査からなる。
\todo{全体の流れの図}

\section{最初の調査}
最初の調査は約86問で、基本情報、性格（TIPI-J）、推し活の中身（12の活動が推し活全体に占める割合）、推し活を通じた気持ちと生活（26項目）、推し活をしたくなるとき（7項目）、WHO-5、主観的幸福感（SHS）、ストレス（PSS-J）からなる。
回答を選ぶと自動で次の質問に進む。

\section{5つの要素と分析画面}
推し活の状態を、推し活の充実、睡眠と生活、お金、義務感と人間関係、SNSと比較の5つの要素で表す。
各要素は該当する項目（7件法）の平均を0〜100点に直したもので、どの要素も高いほど良い向きにそろえる。
60点以上を良好、60点未満を見直したい要素として色分けする。
結果として、幸福度（WHO-5）とストレス度（PSS-J）、推し疲れ（8項目）を示す。
\todo{各要素の項目と式の表（oshifit\_spec.md と画面試作の scoring から）}

\section{群の決め方}
6つの軸のうち5つ以上が60点以上で、かつ WHO-5 の素点が13点以上の人は「今のままで順調」（群A）とし、提案を行わない。
それ以外の人には提案を示し、提案を選んで始めた人を群B、今のまま続けることを選んだ人を群Cとする。

\section{提案の選び方}
提案の候補は15種類で、それぞれが5つの要素または幸福度のいずれかに対応する。
各候補の優先度は、その要素の目安の点数と本人の点数との差、本人の性格との合い方、その活動が推し活に占める比重から計算し、要素が重ならないように上位3つを示す。
\todo{候補の一覧と根拠の確かさの表}

\section{提案の根拠の提示}
提案ごとに、根拠となる研究と、その確かさを3段階（介入研究で確認、実験や近い介入研究がある、調査研究や理論にもとづく）で示す。

\section{2つの未来の提示}
プランA（今のまま）とプランB（選んだ提案を取り入れる）の4週間後の見込みを、本人に似せたアバターと5要素のレーダーチャートで並べる。
見込みは提案ごとの想定の効果を足し合わせた仮の値であり、選んだ提案が多いほど幅を広く示す。
\todo{想定の効果の値の根拠（oshifit\_spec.md 第8節の未決）}

\section{振り返りと最後の調査}
群B・Cの参加者は、4〜5日ごとに計6回の振り返りを行う。
28日目以降に、最初の調査と同じ尺度で最後の調査を行う。

\section{プライバシーと安全への配慮}
名前やメールアドレスは集めず、匿名の参加IDで記録する。
アバター用の写真は端末の中で顔の比率と色を読み取るだけで、保存も送信もしない。
つらい状態が続くときは振り返りを休んでよいことと、相談先を案内する。
